% chapters/13/rangkuman.tex
\bagian{Rangkuman}

\begin{rangkumanbox}
\begin{itemize}
  \item Level pertama disusun secara \istilah{greyboxing}, yaitu dari balok
        abu-abu tanpa aset, agar mekaniknya dapat diuji sebelum aset yang mahal
        dibuat.
  \item Level dua dimensi juga dapat disusun dengan \istilah{tilemap}, yang
        seluruh petaknya cukup memerlukan satu \istilah{Collider}.
  \item Ukuran level ditetapkan lebih dahulu: tinggi lompatan, lebar celah, dan
        tinggi bongkahan. Ketiganya saling bergantung dan menentukan apakah
        levelnya terasa adil.
  \item Benda diam cukup memerlukan \istilah{Collider}; benda bergerak
        memerlukan \istilah{Collider} beserta \istilah{Rigidbody}.
  \item Tiga tipe \istilah{Rigidbody}: \istilah{dynamic} untuk benda yang
        bergerak menurut fisika, \istilah{kinematic} untuk benda yang geraknya
        ditentukan kode, dan \istilah{static} untuk benda diam yang tidak
        memerlukan \istilah{Rigidbody} sama sekali.
  \item \istilah{Trigger} adalah \istilah{Collider} yang tidak menghalangi,
        melainkan hanya memanggil fungsi tertentu pada kode Anda. Ketiga
        syaratnya adalah tanda \istilah{Is Trigger}, keberadaan
        \istilah{Rigidbody}, dan \istilah{layer} yang saling diperiksa.
  \item Kontrol pemain membaca masukan pada \code{Update} dan menggerakkan
        \istilah{Rigidbody} pada \code{FixedUpdate}. Perbedaan keduanya sama
        dengan pembedaan langkah berubah dan langkah tetap pada Bagian~2.2.
  \item Pemeriksaan tanah dilakukan dengan memeriksa lingkaran kecil di bawah
        kaki pemain, bukan dengan membandingkan letak lantai, agar bekerja pada
        segala bentuk permukaan.
  \item Empat masalah fisika yang paling sering terjadi: benda bergetar,
        menembus \istilah{Collider}, jatuh tanpa henti, dan tersangkut pada
        dinding. Seluruhnya memiliki gejala dan penyebab yang khas.
  \item Urutan penelusuran masalah fisika: periksa komponennya lebih dahulu,
        tampilkan bentuk batas benda, periksa pengaturan fisika, baru periksa
        kode Anda. Pada Unity, sebagian besar masalah justru berada pada
        pengaturan komponen.
  \item Kamera dipindahkan pada \code{LateUpdate}, yaitu setelah seluruh
        perhitungan bingkai selesai, dan kelambatannya merupakan salah satu
        pengatur rasa gerak yang paling berpengaruh.
  \item Bab ini menutup \CPMK{3}. Level yang dihasilkan menjadi pijakan bagi
        Bab~XIV sampai Bab~XVI, yang melengkapinya dengan aset, audio, dan
        pengujian.
\end{itemize}
\end{rangkumanbox}
